% part: intuitionistic-logic
% chapter: propositions-as-types
% section: proof-terms

\documentclass[../../../include/open-logic-section]{subfiles}

\begin{document}

\olfileid{int}{pty}{ter}

\olsection{ثبوت ترمونه}

په \Log{N1i} کښې فرضونه په متغيرونو نښه کوو (مثلاً
$!A^x$)، څو د فرض ايستلو نښې وي۔ په \Log{N2i} کښې هر
سېکوېنټ $\Gamma \Sequent !A$ په کيڼ سياق~$\Gamma$
کښې د متغير--!!{formula} جوړو لړ لري۔ نو په \Log{N2i}
کښې د !!a{proof} هر سېکوېنټ يوازې دا نۀ ښيي چې !!{proof}
تر اوسه څه !!{prove}d کړي (يعنې $!A$)، بلکې دا هم ښيي
چې په هر ځای کښې کوم فرضونه پرانيستي دي (يعنې
$\Gamma$)۔ که په هر سېکوېنټ کښې دا معلومات هم شامل
کړو چې $!A$ \emph{څنګه} !!{prove}d شوے، نو څه به وشي؟
د دې دپاره يو \emph{په ترمونو نښه شوے} نظام \Log{tN2i}
تعريفولے شو، چې سېکوېنټونه ئې د $\Gamma \Sequent N : !A$
په بڼه وي۔ $\Gamma$~او $!A$ هماغه پخواني دي: د
$x: !B$ په بڼه جوړو سېټ، او هغه !!{formula} چې په
دغه سېکوېنټ پاې ته رسېدونکي فرعي-!!{proof} ئې له
فرضونو~$\Gamma$ څخه پيروي ښودلې ده۔ $N$~به داسې
ترم وي چې په دغه سېکوېنټ پاې ته رسېدونکے !!{proof}
کوډ کوي۔

هر سېکوېنټ ته ټاکل شوي ترمونه له متغيرونو $x$، $y$،
\dots، او داسې تابع نښو جوړېږي چې تر اوسه ترسره شوي
استنتاجونه کوډ کوي۔ د هرې استنتاجي قاعدې دپاره بېله
تابع نښه پکار ده۔ د !!{introduction} له قواعدو سره
تړلې نښې «جوړوونکې» او د !!{elimination} له قواعدو
سره تړلې نښې «ماتوونکې» بلل کېږي۔ هره تابع نښه د
قاعدې د مقدمو په شمېر آرګومېنټونه اخلي۔ مثلاً د
\Intro{\land} او \Elim{\land}[1] نښې $c^\land$
(دوه‌ځايه) او $d^\land_1$ (يوځايه) کېدے شي۔ په
\Log{tN2i} کښې دغه قواعد داسې کېږي:
\[
\Axiom$\Gamma_1 \fCenter N: !A$
\Axiom$\Gamma_2 \fCenter M: !B$
\RightLabel{\Intro{\land}}
\BinaryInf$\Gamma_1, \Gamma_2 \fCenter c^\land(N, M) : !A \land !B $
\DisplayProof
\quad
\Axiom$\Gamma \fCenter N: !A \land !B$
\RightLabel{\Elim{\land}[1]}
\UnaryInf$\Gamma \fCenter d^\land_1(N) : !A$
\DisplayProof
\]
که استنتاجي قاعده يو فرض باسي، يعنې په يوه مقدمه کښې
نښه لرونکے !!{formula}~$x:!A$ وي چې د پايلې په سياق
کښې نشته، نو په ثبوت ترم کښې هم دغه خبره ښودل پکار
دي۔ له دغو قواعدو سره تړلې تابع نښې اړوند متغيرونه
\emph{تړي}۔ مثلاً د \Intro{\lif} قاعده له مقدمې
$x:!A, \Gamma \Sequent !B$ څخه پايلې
$\Gamma \Sequent !A \lif !B$ ته ځي۔ جوړوونکے
$c^{\lif}$ يو آرګومېنټ اخلي، خو دا هم بايد وښيي چې
متغير $x$ تړل کېږي۔ ثبوت ترم په دې توګه ښيي چې د
نښې~$x$ فرض ايستل شوے دے۔ مثلاً که د مقدمې ثبوت
ترم~$N$ وي، د پايلې ثبوت ترم $c^{\lif}([x]N)$
او قاعده داسې ليکلے شو:
\[
\Axiom$x:!A, \Gamma \fCenter N: !B$
\RightLabel{\Intro{\lif}}
\UnaryInf$\Gamma \fCenter c^{\lif}([x]N) : !A \lif !B$
\DisplayProof
\]
\textbf{د ژباړې د نسخې سمون (OLCMP-208):} د او د
جوړوونکي دويم آرګومېنټ د هماغې دويمې مقدمې لوے
توري ترم شو؛ د سرچينې کوچنے تورى دلته بې‌تعريفه و۔

د !!a{proof} د بشپړې بيا جوړونې دپاره نور معلومات
هم پکار دي۔ که د قاعدې د مقدمو !!{formula}s د پايلې
!!{formula} په بشپړ ډول نۀ ټاکي، دغه معلومات بايد
ترم ته ورزيات کړو۔ په \Intro{\lif} کښې همداسې دي،
نو $c^{\lif}([x:!A]N)$ ليکو۔ په \Intro{\lor} او
~\FalseInt کښې همداسې دي، نو داسې تابع نښې کاروو
چې دغه معلومات ولري؛ مثلاً:
\[
\Axiom$\Gamma \fCenter N:!A$
\RightLabel{\Intro{\lor}[1]}
\UnaryInf$\Gamma \fCenter c^{\lor{!B}}_1(N):!A \lor !B$
\DisplayProof
\quad
\Axiom$\Gamma \fCenter N:\lfalse$
\RightLabel{\FalseInt}
\UnaryInf$\Gamma \fCenter d^\lfalse_{!A}(N):!A$
\DisplayProof
\]
\textbf{د ژباړې د نسخې سمون (OLCMP-209):} د يا د
جوړوونکي ترم ته د مقدمې ترم آرګومېنټ ورګډ شو؛ په
سرچينه کښې له وروستۍ نښې دغه لازم آرګومېنټ ورک و۔

په دغو طريقو داسې سېکوېنټي فطري استنتاجي نظام
جوړولے شو چې هر سېکوېنټ ئې د $\Gamma \Sequent N : !A$
په بڼه وي او $N$~يو \emph{ثبوت ترم} وي۔

دغه ثبوت ترمونه له هغه حساب د ترمونو سره نږدې
اړيکه لري چې \emph{ټايپ‌لرونکے لامبډا حساب} بلل
کېږي۔ د دې اړيکې د ساتلو دپاره به د لامبډا حساب
په ادبياتو کښې کارېدونکې نښې وکاروو، د پورته
عمومي جوړوونکو او ماتوونکو نښو پر ځاے۔ مثلاً د
$c^{\lif}([x:!A]N)$ پر ځاے $\lambd[\typeof{x}{!A}][N]$،
د $d^{\lif}(N,M)$ پر ځاے $(NM)$، د $c^\land(N,M)$
پر ځاے $\pair{N}{M}$، او د~$d^\land_i(N)$ پر ځاے
$\proj{i}{N}$ ليکو۔

\begin{defn}
  ثبوت ترمونه په استقرايي ډول د لاندې قواعدو په وسيله جوړېږي:
  \begin{enumerate}
  \item هر متغير $x$ ثبوت ترم دے (بديهيتونه)۔
  \item که $x$ متغير، $!A$، !!a{formula} او $N$ ثبوت ترم
    وي، نو $\lambd[\typeof{x}{!A}][N]$ هم ثبوت ترم دے~(\Intro\lif)۔
  \item که $N$ او $M$ ثبوت ترمونه وي، نو $(NM)$ هم ثبوت
    ترم دے~(\Elim\lif)۔
  \item که $N$ او $M$ ثبوت ترمونه وي، نو $\pair{N}{M}$
    ثبوت ترم دے~(\Intro\land)۔
  \item که $N$ ثبوت ترم وي، نو $\proj{i}{N}$ هم ثبوت
    ترم دے؛ دلته $i$، $1$ يا~$2$ دے~(\Elim\land[i])۔
  \item که $N$ ثبوت ترم او $!A$ يو فورمول وي، نو
    $\inj[!A]{i}{N}$ ثبوت ترم دے؛ دلته $i$، $1$ يا~$2$
    دے~(\Intro\lor[i])۔
  \item که $M$، $N_1$ او $N_2$ ثبوت ترمونه، او $x_1$ او
    $x_2$ متغيرونه وي، نو $\dcase{M}{x_1}{N_1}{x_2}{N_2}$
    ثبوت ترم دے~(\Elim\lor)۔
  \item که $N$ ثبوت ترم او $!A$، !!a{formula} وي، نو
    $\abort{!A}{N}$ ثبوت ترم دے~(\FalseInt)۔
  \end{enumerate}
  \textbf{د ژباړې د نسخې سمون (OLCMP-210):} د يا د
  جوړوونکي په ماده کښې هماغه فرض شوے ثبوت ترم
  آرګومېنټ شو؛ د سرچينې بل، بې‌تعريفه او په فورمولي
  نښه ليکل شوے ترم سم نۀ و۔
\end{defn}

هر ثبوت ترم د !!a{proof} ترم نۀ دے۔ مثلاً ترم
$\pair{N}{M}$ داسې ثبوت کوډ کوي چې په يوه~\Intro{\land}
استنتاج پاے ته رسېږي، نو پاې-!!{formula} ئې د او
تړنې فورمول~$!A \land !B$ دے۔ دا د يوې~\Elim{\lif}
قاعدې لويه مقدمه نۀ شي کېدے، نو $(\pair{N}{M}P)$ د
صحيح ثبوت ترم نۀ شي کېدے۔ يوازې هغه ثبوت ترمونه
\emph{صحيح} دي چې د \Log{tN2ip} قواعد تعقيبوي۔ دغه
قواعد په \olref{tab:tN2ip} کښې ورکړل شوي دي۔

\subfile{rules-tN2}

\end{document}
